\documentclass[11pt]{article}
\usepackage{amsmath,amssymb,booktabs,geometry}
\geometry{margin=1in}

\title{Step 201: Burnol \(\kappa\) Operationalization}
\author{six-birds-foundations-iii}
\date{}

\begin{document}
\maketitle

\section*{Verdict}

\[
\boxed{\texttt{V\_kappa\_inherited\_branch\_AB\_attempted}.}
\]

The manager-led paper-fetch audit of 2026-05-16 supplies the external Burnol
records that were missing in Steps 177--180 and 191--192.  The old statement
``stuck at \(\kappa\)'' is lifted at the \(\kappa\)-level.  Branches A and B
do not close in this step; instead, new downstream residuals are exposed.

\section*{Paper-Grounded Inherited Records}

The following records are inherited from the manager-led audit, not from
training memory.

\begin{enumerate}
\item Burnol 2004 (math/0112254), Definition 6.1:
\[
  L_a^\Gamma \equiv K_\lambda
  =
  L^2((\lambda,\infty),dt)
  \cap
  F_+\bigl(L^2((\lambda,\infty),dt)\bigr),
  \qquad a=\lambda>0.
\]

\item Burnol 2004, Theorem 6.10:
\[
  A_\infty(s)=\pi^{-s/2}\Gamma(s/2),\qquad
  M(f)(s)=\pi^{-s/2}\Gamma(s/2)\widehat f(s),
\]
where
\[
  \widehat f(s)=\int_0^\infty f(t)t^{-s}\,dt.
\]

\item Burnol 2002 (math/0208121), Theorem 4:
\[
\pi_\lambda(f)
=f-\left[(1-D_\lambda)^{-1}P_\lambda(f)-F_\lambda F_+(f)\right]
-F_+\left[(1-D_\lambda)^{-1}P_\lambda F_+(f)-F_\lambda(f)\right],
\]
where \(P_\lambda\) projects onto \(L^2(-\lambda,\lambda)\),
\[
  F_\lambda=P_\lambda F_+P_\lambda,\qquad D_\lambda=F_\lambda^2.
\]
The manager audit records that \(D_\lambda\) may be represented by the
Dirichlet sinc kernel
\[
  \frac{\sin(2\pi\lambda(x-y))}{\pi(x-y)}
\]
on the even subspace.

\item Burnol 2002, Corollary 5:
\[
f=F_+f\Longrightarrow
\pi_\lambda(f)=f-(1+F_+)(1+F_\lambda)^{-1}P_\lambda(f),
\]
\[
f=-F_+f\Longrightarrow
\pi_\lambda(f)=f-(1-F_+)(1-F_\lambda)^{-1}P_\lambda(f).
\]

\item Burnol 2002, Theorem 8:
\[
  E_\lambda(w)=
  \pi^{-w/2}\Gamma(w/2)
  \left[
    \lambda^{1/2-w}
    +
    \frac{\sqrt\lambda}{2}
    \int_\lambda^\infty
      \bigl(\psi_+^\lambda(t)-\psi_-^\lambda(t)\bigr)t^{-w}\,dt
  \right],
\]
with
\[
\psi_+^\lambda(t)
 =
2\cos(2\pi\lambda t)
-
F_+\left[
(1+F_\lambda)^{-1}P_\lambda(2\cos(2\pi\lambda y))
\right](t),
\]
\[
\psi_-^\lambda(t)
 =
2\cos(2\pi\lambda t)
+
F_+\left[
(1-F_\lambda)^{-1}P_\lambda(2\cos(2\pi\lambda y))
\right](t).
\]

\item Burnol 2002, equation 1:
\[
K_a^\Gamma(z,w)
\equiv
K_{B(E_\lambda)}(z,w)
=
\frac{
E_\lambda(z)\overline{E_\lambda(w)}
-
E_\lambda(1-z)\overline{E_\lambda(1-w)}
}{z+w-1}.
\]
\end{enumerate}

\section*{Derivation of \(\kappa\)}

Step 152--153 gives
\[
  T_a=M_\Gamma J_aU_\infty^{-1}.
\]
Therefore
\[
  T_a^*=(U_\infty^{-1})^*J_a^*M_\Gamma^*.
\]
With \(K_a^\Gamma\) now explicit, the pulled evaluator is
\[
  \kappa_{a,w,k}(\tau)
  =
  \left(
  T_a^*\partial_{\overline w}^{\,k}K_a^\Gamma(\cdot,w)
  \right)(1/2+i\tau).
\]
For the Step 164 finite carrier \(a=\lambda=1/2\), \(k=0\), and
\(\rho_i\in\{\rho_1,\rho_2,\rho_3\}\),
\[
  \kappa_i(\tau)=
  \left(
  T_{1/2}^*K_{1/2}^\Gamma(\cdot,\rho_i)
  \right)(1/2+i\tau).
\]

\section*{Branch B SL164.1}

Because \(T_a\) is the inherited unitary transport, the Gram entries are now
symbolic closed forms:
\[
  G_{ij}=\langle \kappa_i,\kappa_j\rangle
  =
  K_{1/2}^\Gamma(\rho_j,\rho_i)
\]
up to the inherited convention for linearity of the inner product.  Thus
\[
  e_i=\frac{\kappa_i}{\sqrt{G_{ii}}}.
\]
The finite commutator matrix is
\[
  c_{ij}(\ell)
  =
  \left\langle
    e_i,(I-P_\infty)M_{m_\ell}P_\infty e_j
  \right\rangle,
  \qquad
  m_\ell(s)=e^{-\ell(1/2-s)}.
\]
Using Step 173's \(K_\infty^{op}\), this becomes an explicit double
operator integral in the Mellin boundary variable.  The finite residual
\(\Xi_{\mathrm{matrix\_source}}\) is therefore no longer blocked by
missing \(\kappa\), but it is not numerically decided here.  The new
sub-residual is certified evaluation of \(E_{1/2}\), the
\((1\pm F_\lambda)^{-1}\) resolvents, and the \(K_\infty^{op}\) commutator
integrals.

\section*{Branch A G2--G5}

For \(H_\eta=\overline{\mathrm{span}\{\eta^a_{\rho,k}\}}\), each evaluator
has the explicit transported form above.  Matrix entries of the active
Calkin object are
\[
  \langle \eta_i,C_\ell P_\eta\eta_j\rangle
  =
  \langle
  \kappa_i,(I-P_\infty)M_{m_\ell}P_\infty\kappa_j
  \rangle.
\]
Hence the matrix-entry attack is no longer blocked by \(\kappa\).
However, Branch A closure requires a completed Calkin conclusion:
\[
  C_\ell P_\eta\in K_\eta
  \quad\text{or}\quad
  \|q_\eta(C_\ell P_\eta)\|>0,
\]
or the original G2--G5 bridge theorem.  The explicit formula exposes the
new tasks: summability of the infinite matrix, a compactness theorem, a
positive essential-norm Weyl sequence, or a faithful boundary-symbol theorem.

\section*{CTMT Re-Evaluation}

Branch A and Branch B are no longer CTMT-stuck at \(\kappa\).  Their
\(\kappa\)-level terminality is lifted.  CTMT itself remains valid: it now
records a resolved terminality instance that produced deeper sub-residuals.
Branch C remains CTMT-foreclosed-numerical and is unchanged by the new
Burnol records.

\section*{Nonclaim}

This step does not prove RH, close \(\Xi_{BC}\), close Branch A, close Branch
B, or weaken Branch C's numerical foreclosure.  It advances Path 1 from
``blocked at \(\kappa\)'' to ``\(\kappa\) available; downstream gates open.''

\end{document}

